% The optional title feeds the contents and running heads; the line
% break balances the two lines of the opening.
\chapter[Heatwave Response of Shoots and Sediment]{Heatwave Response\\ of Shoots and Sediment}
\label{ch:heatwave-response}

\section{A threshold in shoot density}
\label{sec:threshold}

Across the six summers from 2019 to 2024, the cumulative heat dose of
\cref{eq:heat-dose} ranged from zero in the cool summer of 2021 to more
than thirty degree-heating-days in the shallow meadows of 2022. To test
whether shoot density responds gradually or abruptly, I compared a linear
model with the segmented model of \textcite{vandermeer2016segmented},
\begin{equation}
  \log \frac{N_{\text{Oct}}}{N_{\text{Jun}}} =
  \begin{cases}
    \alpha + \beta_1 H, & H \le \tau,\\[0.4ex]
    \alpha + \beta_1 \tau + \beta_2 (H - \tau), & H > \tau,
  \end{cases}
  \label{eq:segmented}
\end{equation}
where \(N\) is the mean shoot density of a meadow and \(\tau\) the
breakpoint. The segmented model was strongly preferred
(\(\Delta\mathrm{AIC} = 23.4\)), with a breakpoint at
\(\tau = 14.2\) degree-heating-days (95\,\% interval 11.8--16.9).

\begin{figure}[tb]
  \centering
  \begin{tikzpicture}
    \begin{axis}[
        width=0.92\linewidth, height=0.5\linewidth,
        xmin=0, xmax=34, ymin=-2.5, ymax=0.4,
        xlabel={Cumulative heat dose \(H\) (degree-heating-days)},
        ylabel={\(\log(N_{\text{Oct}}/N_{\text{Jun}})\)},
        xtick={0,5,...,30}, ytick={-2,-1.5,-1,-0.5,0},
        axis x line*=bottom, axis y line*=left,
        clip=false,
      ]
      \fill[accent!8] (axis cs:11.8,-2.5) rectangle (axis cs:16.9,0.4);
      \addplot[only marks, mark=*, mark size=1.3pt, black!55] coordinates {
        (0.0,0.12) (0.4,0.05) (1.2,0.09) (2.1,-0.02) (2.8,0.04)
        (3.5,-0.06) (4.4,0.01) (5.1,-0.08) (6.0,-0.03) (6.8,-0.12)
        (7.5,-0.07) (8.3,-0.15) (9.1,-0.10) (9.9,-0.18) (10.6,-0.13)
        (11.4,-0.22) (12.3,-0.19) (13.0,-0.27) (13.8,-0.24) (14.6,-0.33)
        (15.5,-0.40) (16.4,-0.58) (17.2,-0.60) (18.5,-0.82) (19.6,-0.86)
        (20.8,-1.08) (22.1,-1.15) (23.5,-1.39) (25.0,-1.48) (26.7,-1.76)
        (28.3,-1.86) (30.4,-2.18) (32.1,-2.31)
      };
      \addplot[accent, line width=1pt] coordinates {
        (0,0.08) (14.2,-0.28) (32.5,-2.39)
      };
      \draw[accent, dashed, line width=0.4pt]
        (axis cs:14.2,-2.5) -- (axis cs:14.2,0.4);
      \node[accent, anchor=south west, font=\footnotesize]
        at (axis cs:14.6,0.15) {\(\tau = 14.2\)};
    \end{axis}
  \end{tikzpicture}
  \caption{Change in shoot density over the summer against cumulative
    heat dose, for all meadows and years. The line is the segmented fit of
    \cref{eq:segmented}; the shaded band is the 95\,\% interval of the
    breakpoint.}
  \label{fig:threshold}
\end{figure}

\Cref{fig:threshold} shows the fit. Below the threshold the decline is
slight, about two and a half per cent of shoots for each additional
degree-heating-day, and is offset by recruitment in the following spring.
Above it, the slope steepens almost fivefold. The threshold lies close to
the dose at which \textcite{okafor2018thermal} observed the collapse of
net photosynthesis in laboratory cultures, which suggests a physiological
rather than a merely statistical boundary.\footnote{The chamber
  experiment of \cref{sec:censuses} reproduced the breakpoint within its
  confidence interval in two of three years; in the third, a storm
  removed the chambers before the heat dose had accumulated.}

\section{Loss of sediment carbon}
\label{sec:carbon-loss}

Cores taken one year after a summer above the threshold contained less
organic carbon in their upper ten centimetres than cores from the same
meadow before the event (\cref{tab:carbon}). Below ten centimetres no
change could be detected, consistent with the depth to which oxygen
penetrates once the rhizome mat has thinned
\parencite{aagaard2021remineralisation}.

\begin{table}[b]
  \centering\small
  \caption{Organic carbon in the upper 10 cm of sediment before and one
    year after summers above the heat-dose threshold (mean of three cores
    per meadow).}
  \label{tab:carbon}
  \begin{tabular*}{\linewidth}{@{\extracolsep{\fill}}l r r r@{}}
    \toprule
    & \multicolumn{2}{c}{C stock (Mg ha\(^{-1}\))} & \\
    \cmidrule(lr){2-3}
    Meadows & Before & After & Loss (\%) \\
    \midrule
    1--3   & 16.2 & 11.2 & 30.9 \\
    4--6   & 19.8 & 14.6 & 26.3 \\
    7--9   & 22.4 & 18.4 & 17.9 \\
    \bottomrule
  \end{tabular*}
\end{table}

Recovery was slow. \Cref{fig:recovery} follows the three shallowest
meadows after the 2022 heatwave alongside a meadow that stayed below the
threshold; shoot density returned to its pre-event level within one
season in the latter, but had regained only two thirds of it after two
years in the former, in line with observations from the open North Sea
\parencite{sato2023recovery}.

\begin{figure}[tb]
  \centering
  \begin{tikzpicture}
    \begin{axis}[
        width=0.92\linewidth, height=0.44\linewidth,
        xmin=-0.2, xmax=4.2, ymin=0, ymax=120,
        xlabel={Years after the 2022 heatwave},
        ylabel={Shoot density (\% of 2021)},
        xtick={0,1,2,3,4}, ytick={0,25,50,75,100},
        axis x line*=bottom, axis y line*=left,
        legend pos=south east,
      ]
      \addplot[accent, line width=1pt, mark=*, mark size=1.2pt]
        coordinates {(0,38) (1,49) (2,66) (3,81) (4,94)};
      \addlegendentry{Meadows 1--3, above threshold}
      \addplot[black!55, line width=0.8pt, mark=o, mark size=1.4pt]
        coordinates {(0,86) (1,101) (2,99) (3,103) (4,100)};
      \addlegendentry{Meadow 11, below threshold}
    \end{axis}
  \end{tikzpicture}
  \caption{Recovery of shoot density after the 2022 heatwave in meadows
    above and below the threshold.}
  \label{fig:recovery}
\end{figure}

\section{Implications for carbon accounting}
\label{sec:accounting}

Scaling the losses in \cref{tab:carbon} by meadow area and by the
frequency with which each meadow crossed the threshold between 2019 and
2024 gives an average return of about \(0.6\) megagrams of carbon per
hectare per year. Set against a mean burial rate of \(2.4\) megagrams
\parencite{holmgren2019burial}, this means that current guidelines, which
treat buried carbon as permanent \parencite{bcwg2023guidelines},
overestimate the net benefit of the estuary's meadows by about a quarter.
The summer heat dose, cheap to measure and available weeks before shoots
decline, offers managers a direct way to anticipate those losses.
